% Capitolo 4: Campioni

I campioni sono stati scelti per coprire i tre fenomeni polarimetrici di interesse (birifrangenza, attività ottica e fotoelasticità), scegliendo materiali trasparenti compatibili con la configurazione in trasmissione. Alcuni possiedono parametri nominali noti, utili per una validazione quantitativa; altri offrono strutture spaziali che mettono alla prova la capacità di imaging del sistema. La \cref{tab:campioni} ne riassume le caratteristiche principali; tutti sono inseriti nel porta-campione sulla rotaia, tra sorgente e analizzatore, e le superfici vengono pulite prima dell'acquisizione.

\begin{table}[H]
    \centering
    \caption{Campioni analizzati con il polarimetro a immagine.}
    \label{tab:campioni}
    \small
    \begin{tabularx}{\textwidth}{@{}l >{\raggedright\arraybackslash}X >{\raggedright\arraybackslash}X@{}}
        \hline
        \textbf{Campione} & \textbf{Fenomeno} & \textbf{Validazione attesa} \\
        \hline
        Lamina $\lambda/2$  & Birifrangenza & $\delta \approx 180^\circ$, AoLP ruotato \\
        Lamina $\lambda/4$ & Birifrangenza & $|S_3/S_0|$ massimo \\
        Nastro adesivo (1--5 strati) & Birifrangenza incrementale & $\delta \propto N$ \\
        Soluzione di saccarosio & Attività ottica & $\alpha_\text{rot} \propto 1/\lambda^2$ \\
        Righello in plastica & Fotoelasticità (stress residuo) & Frange isocromatiche \\
        Trave a sbalzo & Fotoelasticità (stress applicato) & Birifrangenza indotta dal carico \\ % TODO: check (cella modificata da AI 2026-05-21)
        \hline
    \end{tabularx}
\end{table}


\section{Campioni birifrangenti}
\label{sec:campioni-birifrangenti}

Le lamine d'onda commerciali ($\lambda/2$ e $\lambda/4$, lunghezza d'onda di design di $\SI{633}{\nano\meter}$, cfr.\ \cref{sec:fenomeni}) sono il primo test di validazione quantitativa: la ritardanza nominale è specificata dal costruttore e la risposta è spazialmente uniforme, il che consente un confronto diretto tra valore misurato e valore atteso. La lamina $\lambda/4$, orientata a $\SI{45}{\degree}$ rispetto alla polarizzazione incidente, produce luce circolarmente polarizzata con $|S_3/S_0| \approx 1$; la lamina $\lambda/2$ ruota il piano di polarizzazione di un angolo pari al doppio della sua orientazione, con ritardanza nominale $\delta \approx 180^\circ$. L'ordine delle due lamine non è dichiarato dal costruttore, sicché a priori se ne conosce soltanto la ritardanza nominale alla lunghezza d'onda di design: l'analisi del \cref{ch:risultati} le caratterizzerà a posteriori, mostrando che la $\lambda/4$ è compatibile con un ritardatore a singolo ordine, mentre la $\lambda/2$ si rivelerà multi-ordine, circostanza che ne fa un caso di studio sull'ambiguità di avvolgimento anziché un semplice riferimento di ritardanza nota.

Il vetrino con strati di nastro adesivo trasparente (da 1 a 5 strati, in configurazione simmetrica $1$-$2$-$3$-$4$-$5$-$4$-$3$-$2$-$1$) è un campione a birifrangenza incrementale controllata: % TODO: check (spiegazione birifrangenza nastro inserita da AI 2026-05-21)
il film polimerico del nastro, stirato uniassialmente in fase di produzione, ha le catene molecolari allineate lungo una direzione preferenziale, e questa anisotropia strutturale gli conferisce una birifrangenza intrinseca. La ritardanza totale cresce approssimativamente come
\begin{equation}
\delta_{\text{tot}} \approx N \cdot \delta_{\text{singolo}},
\label{eq:ritardo-strati}
\end{equation}
dove $N$ è il numero di strati e $\delta_{\text{singolo}}$ il ritardo introdotto da un singolo strato. La simmetria del campione fornisce un controllo interno di consistenza: le regioni con lo stesso numero di strati ai due lati del vetrino devono restituire ritardanze concordanti.


\section{Attività ottica}
\label{sec:campione-attivita-ottica}

Una soluzione di saccarosio (in una boccetta con cammino ottico $l = \SI{13}{\milli\meter}$) consente di verificare due aspetti: la rotazione uniforme del piano di polarizzazione, osservabile come variazione dell'AoLP, e la legge di dispersione $\alpha_\text{rot} \propto 1/\lambda^2$ (modello di Drude), sfruttando i tre canali colore del sensore. Il saccarosio destrogiro ha rotazione specifica $[\alpha]_D = \SI{66.47}{\degree\per(\deci\meter\cdot\gram\per\milli\liter)}$ alla riga D del sodio.


\section{Campioni fotoelastici}
\label{sec:campioni-fotoelastici}

Un righello in plastica trasparente presenta tensioni residue da stampaggio a iniezione che generano birifrangenza non uniforme nel materiale. Quando posto tra polarizzatore e analizzatore, le regioni a diverso stato di tensione producono le frange isocromatiche e isocline tipiche dei materiali fotoelastici (cfr.\ \cref{sec:fenomeni}). L'analisi con il polarimetro a immagine consente di mappare spazialmente la distribuzione delle tensioni interne.

Una trave in materiale trasparente, vincolata a un'estremità, è stata analizzata in due condizioni: scarica, per verificare l'eventuale presenza di tensioni residue, e sotto carico, con un peso applicato all'estremità libera. % TODO: check (paragrafo trave riscritto da AI 2026-05-21: frange -> variazione birifrangenza, qualitativo)
Lo sforzo di flessione varia linearmente lungo lo spessore e inversamente con la distanza dal vincolo, inducendo per effetto fotoelastico una birifrangenza non uniforme lungo la trave. Nelle condizioni di questo esperimento le frange isocromatiche non risultano nettamente risolte: si osserva piuttosto una variazione delle proprietà di birifrangenza, rilevabile soprattutto come modulazione di $S_3$ fra trave carica e scarica. L'analisi su questo campione resta perciò qualitativa.